\documentclass[11pt]{beamer}
\usetheme{metropolis}

\usepackage[french]{babel}
\usepackage[utf8]{inputenc}
\usepackage[T1]{fontenc}
\usepackage{booktabs}
\usepackage{array}

\title{Type de maquette : web ou mobile ?}
\subtitle{Justification du choix retenu pour la construction de l'arbre}
\date{}
\author{}

\begin{document}

\maketitle

\begin{frame}{Sommaire}
  \tableofcontents
\end{frame}

\section{La question posée}

\begin{frame}{Deux options possibles}
  Pour évaluer notre méthode de construction d'arbre, nous avons besoin
  d'un type de maquette source précis.

  \vspace{1em}

  \begin{columns}
    \column{0.5\textwidth}
    \textbf{Mobile}
    \begin{itemize}
      \item Interfaces Android / iOS
      \item Écrans plus simples
      \item Datasets figés existants (RICO, AMP)
    \end{itemize}

    \column{0.5\textwidth}
    \textbf{Web}
    \begin{itemize}
      \item Pages de sites réels
      \item Mises en page riches et imbriquées
      \item Vérité terrain accessible directement (DOM)
    \end{itemize}
  \end{columns}

  \vspace{1em}
  \textbf{Choix retenu : le web.} Les sections suivantes justifient ce choix.
\end{frame}

\section{Argument 1 --- La complexité correspond à notre problème de recherche}

\begin{frame}{Le vide identifié dès le départ}
  Notre point de départ (limite de DeclarUI) :

  \begin{quote}
    Échec sur le \textbf{regroupement d'éléments en sections et listes
    complexes} à l'intérieur d'une page.
  \end{quote}

  \vspace{1em}
  Une interface mobile est en général :
  \begin{itemize}
    \item plus \textbf{plate} (défilement vertical simple)
    \item avec peu de colonnes simultanées
    \item avec moins d'éléments visibles en même temps
  \end{itemize}

  \vspace{1em}
  \alert{Choisir le mobile reviendrait à choisir le cas le plus simple, alors
  que notre problème de recherche vise justement le cas complexe.}
\end{frame}

\begin{frame}{Preuve chiffrée --- DCGen (Divide-and-Conquer)}
  Statistiques mesurées par DCGen sur un jeu de données de sites web réels :

  \vspace{1em}
  \begin{table}
    \centering
    \begin{tabular}{lccc}
      \toprule
      Métrique & Min & Max & Moyenne \\
      \midrule
      Longueur (tokens) & 383 & 868\,108 & 99\,431 \\
      Nombre de balises & 2 & 16\,331 & 961 \\
      Profondeur du DOM & 2 & 40 & 17 \\
      Balises uniques & 2 & 66 & 26 \\
      \bottomrule
    \end{tabular}
  \end{table}

  \vspace{1em}
  \textbf{Une profondeur moyenne de 17 niveaux d'imbrication} : exactement le
  type de structure que notre séparation géométrie/sens doit savoir gérer.
\end{frame}

\section{Argument 2 --- La collecte de vérité terrain}

\begin{frame}{Vérité terrain : un besoin central de l'évaluation}
  Pour mesurer la qualité de notre arbre (Tree Edit Distance, Container
  Match, TreeBLEU\dots), il faut un \textbf{arbre de référence} auquel se
  comparer.

  \vspace{1em}
  \begin{columns}
    \column{0.5\textwidth}
    \textbf{Web}
    \begin{itemize}
      \item Le DOM est directement lisible
      \item Accessible sur des milliers de sites, gratuitement
      \item Aucun outil spécialisé requis
    \end{itemize}

    \column{0.5\textwidth}
    \textbf{Mobile}
    \begin{itemize}
      \item Nécessite la décompilation d'APK
      \item Ou des outils d'accessibilité complexes
      \item Dépendance à des datasets déjà figés
    \end{itemize}
  \end{columns}

  \vspace{1em}
  \alert{Avec un calendrier serré (soutenance fin novembre), la simplicité de
  collecte du web est un critère décisif.}
\end{frame}

\section{Argument 3 --- Cohérence avec les choix déjà faits}

\begin{frame}{Tout notre travail pointe déjà vers le web}
  \begin{itemize}
    \item Notre détecteur retenu (ImageCLEF DrawnUI 2021, tâche Screenshot)
      est entraîné sur \textbf{9\,630 captures de sites web réels}
    \item La limite que nous avions nous-mêmes notée sur Screen Parsing :
      \textit{``entraîné sur mobile (AMP/RICO) --- à vérifier sur des
      maquettes web''}
    \item Les deux méthodes les plus proches de notre approche
      (\textbf{Prototype2Code}, \textbf{DCGen}) travaillent toutes les deux
      sur des captures de sites web
  \end{itemize}

  \vspace{1em}
  \alert{Basculer vers le mobile obligerait à reprendre le détecteur, les
  données de comparaison et une partie de l'état de l'art déjà construit.}
\end{frame}

\section{Synthèse}

\begin{frame}{Synthèse}
  \begin{table}
    \centering
    \begin{tabular}{lp{7.5cm}}
      \toprule
      Argument & Constat \\
      \midrule
      Complexité & Le web correspond au vide de recherche identifié \\
      Vérité terrain & Le DOM est directement accessible, sans outil lourd \\
      Cohérence & Détecteur, état de l'art et travaux proches sont déjà sur le web \\
      \bottomrule
    \end{tabular}
  \end{table}

  \vspace{1.5em}
  \centering
  \Large \textbf{Choix retenu : maquettes web}
\end{frame}

\begin{frame}[standout]
  Merci --- questions et retours bienvenus
\end{frame}

\end{document}
